% Generated draft: requires mathematical and source review.
\documentclass[11pt]{article}
\usepackage[margin=25mm]{geometry}
\usepackage{fontspec}
\setmainfont{Latin Modern Roman}
\usepackage{amsmath,amssymb,mathtools,unicode-math}
\setmathfont{Latin Modern Math}
\usepackage{enumitem,needspace,xurl,hyperref}
\setlength{\parindent}{0pt}
\setlength{\parskip}{5pt}
\setlength{\emergencystretch}{3em}
\newcommand{\N}{\mathbb{N}}
\newcommand{\Z}{\mathbb{Z}}
\newcommand{\Q}{\mathbb{Q}}
\newcommand{\R}{\mathbb{R}}
\newcommand{\C}{\mathbb{C}}
\newcommand{\F}{\mathbb{F}}
\newcommand{\A}{\mathbb{A}}
\newcommand{\PP}{\mathbb{P}}
\newcommand{\E}{\mathbb{E}}
\newcommand{\eps}{\varepsilon}
\newcommand{\dd}{\,\mathrm{d}}
\begin{document}
\section*{UnsolvedMath ID: 30005869}
\textbf{Source problem code:} catalogue entry
\subsection{Definitions and mathematical statement}
Let \(X\) and \((\omega(t))_{t\in I}\) denote the complex manifold and Kähler--Ricci flow intended in the supplied source record, with \(I\) its forward time interval. A Kähler form \(\omega(t)\) determines a Hermitian metric \(g(t)\). Its Ricci form is locally
\[
\operatorname{Ric}(\omega(t))=-\sqrt{-1}\,\partial\bar\partial\log\det(g_{j\bar k}(t)),
\]
and \(|\operatorname{Ric}(\omega(t))|_{\omega(t)}(x)\) denotes its pointwise norm with respect to the evolving metric at \(x\in X\).

The recorded conjectural target is the existence of a finite constant \(C\), independent of space and time, such that
\[
\sup_{t\in I}\sup_{x\in X}
|\operatorname{Ric}(\omega(t))|_{\omega(t)}(x)\le C.
\]
The constant may depend on the fixed geometric data and initial metric; the record requests no uniformity across different initial data. If the intended flow exists for every forward time, then \(I=[0,\infty)\), and the displayed bound is required on that entire interval, including the initial time.

\textbf{Unresolved hypotheses.} The supplied record does not specify the geometric assumptions on \(X\), its dimension, compactness, initial metric, any relevant fibration, the normalization of the flow, or hypotheses guaranteeing existence for all forward time. Consequently, the display states the requested bound for the intended flow but does not determine a complete conjecture quantified over a specified class of manifolds and flows. In particular, no exception near singular fibers is allowed in the requested global bound, should such fibers occur in the intended setting.

The original excerpt also describes a conditional limiting argument involving rescaled flows \(\widetilde\omega_i\), points \(x_i\), and the normalization
\[
|\operatorname{Ric}(\widetilde\omega_i(0))|_{\widetilde\omega_i(0)}(x_i)=1.
\]
It proposes an eternal unnormalized limiting Kähler--Ricci flow, asserts conditionally that a cited Song--Tian result would make that limit scalar-flat and hence Ricci-flat, and expresses hope that the displayed normalization passes to the limit. The rescaling, convergence hypotheses, and cited result are not supplied, and the excerpt ends mid-sentence. These remarks are source motivation, not additional established conclusions or a completed argument.
\subsection{Short English statement}
The question asks whether the Ricci curvature of the intended Kähler--Ricci flow has one finite bound valid at every point and every forward time, measured using the evolving metric. The supplied record omits the geometric hypotheses and flow normalization needed to specify the full conjecture. Its supplied partially solved label and literature assessment require independent review.
\subsection{Sources}
\begin{itemize}
\item Local UnsolvedMath registry, \texttt{data/unsolvedmath/problems.json}, supplied record with permanent ID \(30005869\), exact source problem code \texttt{catalogue entry}, title ``Global Ricci Bounds along Kähler--Ricci Flow.'' The record gives proposed year 2024 and names no proposer.
\item Supplied primary report citation: \emph{Geometrie} (2024), Oberwolfach Reports, pp.\ 1563--1614, DOI \texttt{10.4171/owr/2024/28}, \url{https://doi.org/10.4171/owr/2024/28}. The supplied background describes Oberwolfach Reports as documenting mathematical workshops and their open-problem discussions. No specific page or contribution locating this conjecture is supplied.
\item Supplied literature reference: ``Fong--Lee higher-order Kähler--Ricci estimates,'' \url{https://doi.org/10.1016/j.jfa.2021.109235}.
\item Supplied literature reference: ``Collapsing immortal Kähler--Ricci flows,'' \url{https://doi.org/10.1017/fmp.2025.10}.
\item Supplied literature reference: ``original OWR report,'' \url{https://doi.org/10.4171/owr/2024/28}.
\item Literature assessment dated exactly \texttt{2026-08-22}: the supplied curation labels the problem partially solved and reports uniform Ricci bounds on compact subsets away from singular fibers in important cases, including biholomorphic generic fibers, while describing the global bound as open in general. This is supplied dated curation, not an independently verified assessment. The source explicitly requests independent verification of time-sensitive status and cited claims.
\item The supplied record labels its statement \texttt{corrected\_verified} and says it was checked against the official report and preserves its hypotheses, notation, and target. This is supplied verification metadata, not verification performed for this draft; the hypotheses are not reproduced in the supplied excerpt. The record's update timestamp is \texttt{2026-09-16T12:06:53Z}.
\item The excerpt attributes a preceding result to Fong--Lee, reference [3], developed from an elliptic PDE result [7] using techniques beyond Yau's estimates, and invokes Song--Tian [11] for its conditional scalar-curvature discussion. The corresponding numbered bibliography entries and precise theorem hypotheses are not supplied.
\end{itemize}
\end{document}
